\section{Related Work}
\label{sec:related}

\subsection{FPGA Prototyping Infrastructure}

FASED separates a target DRAM model from its FPGA host implementation and makes
model timing parameters configurable without rebuilding the complete target
\cite{biancolin2019fased}. LEAP provides reusable services and latency-
insensitive interfaces for FPGA applications \cite{fleming2014leap}. RIFFA
offers a portable host--FPGA communication interface \cite{jacobsen2013riffa}.
These systems show the value of stable FPGA infrastructure and host/target
decoupling. \sys applies that principle to a different object: a real
peripheral subsystem that a processor DUT discovers and operates with its
native device driver. Its challenge includes the coupled configuration, BAR,
queue, DMA, completion, and interrupt contract.

\subsection{Modular and Networked FPGA Systems}

Galapagos abstracts heterogeneous CPU/FPGA nodes and makes communication layers
replaceable without changing the application \cite{eskandari2019galapagos}.
Network FPGA clusters similarly separate logical accelerator groupings from
physical datacenter placement \cite{tarafdar2017clusters}. These works organize
compute kernels and their communication. \sys instead organizes devices
presented to an FPGA-hosted processor: the endpoint must participate in PCIe
discovery and preserve a vendor protocol expected by an unmodified OS driver.

\subsection{FPGA Operating Systems and Device Access}

Feniks provides an FPGA OS layer, abstract I/O interfaces, direct PCIe access to
server resources, and a design for remote resource access
\cite{zhang2017feniks}. HybridOS integrates CPU and reconfigurable accelerators
under OS runtime support \cite{kelm2008hybridos}. Their service target is an
FPGA application or accelerator. \sys serves a processor DUT and presents a
driver-compatible device hierarchy. It does not claim remote access or PCIe
peer-to-peer transfer as novel in isolation. Its contribution is the software-
defined endpoint view plus end-to-end mediation needed to take the DUT from
enumeration through native I/O.

\subsection{PCIe Virtualization, Disaggregation, and Composable Fabrics}

Device passthrough and composable PCIe fabrics can assign physical functions or
fabric-attached resources to hosts, and PCIe switches can route those functions
between hosts and devices. Their control plane operates on real PCIe functions
and links; the primary abstraction is resource assignment in the fabric.
\sys instead constructs a logical configuration hierarchy at an FPGA prototype
boundary and may adapt commands and DMA addresses before a physical backend
executes them. Its abstraction is therefore different in two ways: the DUT
sees a software-selected device presentation, and the host coordinator
preserves device-specific control/data semantics across the boundary. \sys is
complementary to composable PCIe hardware, but its evaluation target is
system-software validation rather than transparent extension of the physical
PCIe link.

\subsection{Position of SCOPE}

Across these systems, the reusable idea is a stable interface between a
workload and a changing implementation. SCOPE applies that idea to a processor
DUT that must run a native peripheral driver. The contribution is not a new
PCIe transport, a new NVMe/NIC protocol, or a claim that remote DMA is novel in
isolation. It is the method of combining device-presentation--backend-
implementation decoupling with control-semantics--data-path decoupling, then
constraining the split with configuration--route, submission--DMA, and
data--completion coherence.
